\subsubsection{Greedy Search}
Greedy Search is the basic pursuit strategy used in the \textit{simple} level. 
At the beginning of each round, the system runs a Breadth-First Search (BFS) 
from the current player position to construct a distance field. This distance 
field stores the real grid-step distance from every reachable cell to the player, 
while considering walls and wrap-around movement.

After the distance field is computed, each monster checks all cells that it can 
move to and selects the cell with the shortest distance to the player. This 
process is repeated in every round. The system also prevents two monsters from 
occupying the same cell. If a monster is stuck or no valid BFS path to the player 
exists, Manhattan distance is used as a fallback heuristic, so that the monster 
can still make forward progress.

Greedy Search is simple and efficient. In each round, the system only needs to 
perform one BFS computation and a constant number of comparisons for each monster. 
However, its main limitation is that the monster behaves passively. It only reacts 
to the player's current position and cannot predict where the player may move next. 
Against a clever player, the monster may be trapped or forced into inefficient 
chasing patterns, which increases the time needed to catch the player. Therefore, 
Greedy Search serves as a well-known baseline strategy and provides the foundation 
for more advanced pursuit algorithms.